% Appendix A -- Godkendt projektbeskrivelse

\chapter{Godkendt projektbeskrivelse}
\label{AppendixA}

Den projektbeskrivelse, der ligger til grund for nærværende
rapport, er gengivet nedenfor i den ordlyd, som blev godkendt ved
projektets start (09.\,04.\,2026).

\section*{Begrundelse for valg af emne}

Dette projekt er valgt med udgangspunkt i en konkret problemstilling
stillet af Teknologisk Institut, hvor der arbejdes med udvikling og
implementering af en humanoid robot baseret på eksisterende design
og software.

Projektet har høj faglig relevans, da det kombinerer flere centrale
discipliner inden for automationsteknologi, herunder mekanisk
konstruktion, elektrisk integration samt styring af aktuatorer.

\section*{Problemformulering}

Hvordan kan en humanoid robot, baseret på eksisterende 3D-design og
software, fremstilles, samles og integreres til en funktionel enhed,
hvor mekaniske og elektriske systemer korrekt understøtter robottens
bevægelser?

Projektet omfatter:

\begin{itemize}
    \item 3D-print og efterbehandling af robottens komponenter
    \item Mekanisk samling af robottens delsystemer
    \item Elektrisk opbygning og korrekt wiring af aktuatorer og
          styresystem
    \item Integration og afprøvning af eksisterende software
\end{itemize}

Der arbejdes med at sikre korrekt samspil mellem systemets
delkomponenter samt mulighed for tilpasning og optimering af
udvalgte komponenter og styring med henblik på at forbedre
funktionalitet og stabilitet.

\clearpage
\section*{Afgrænsning}

Projektet afgrænses til at fokusere på udvikling, fremstilling og
integration af en humanoid robotarm som et selvstændigt delsystem.

Robotarmen er valgt som fokusområde, da den indeholder centrale
mekaniske og styringsmæssige udfordringer, herunder bevægelse i
flere led, belastning af komponenter samt koordinering af aktuatorer.

Projektet vil derfor ikke nødvendigvis omfatte en fuld
færdiggørelse af hele den humanoide robot, men der arbejdes med en
modulær tilgang.

\section*{Målgruppe}

Projektet henvender sig primært til teknikere og fagpersoner inden
for automation og robotteknologi, som arbejder med opbygning,
integration og idriftsættelse af robotsystemer.

Fokus er på en praktisk og teknisk formidling.
